
\chapter{Evaluation}
\label{chap:evaluation}


% \ifpdf
%     \graphicspath{{7_conclusions/figures/PNG/}{7_conclusions/figures/PDF/}{7_conclusions/figures/}}
% \else
%     \graphicspath{{7_conclusions/figures/EPS/}{7_conclusions/figures/}}
% \fi

{\color{red} This chapter restates the research question, synthesises the findings from all three experiments, and provides a direct answer supported by the evidence gathered across the three experiments.}


\section{Research Question}


The thesis was organized around a single research question:

\begin{quote}
\textit{``How does the performance of value-based Deep Q-Network agents
compare to actor-critic agents (specifically PPO and A2C) in
discrete-action environments, including their sample efficiency and
training consistency?''}
\end{quote}

Three experiments were designed to answer this question progressively:
Experiment~1 on classic-control environments with MLP policies,
Experiment~2 on Atari with CNN policies at a 5M-step budget, and
Experiment~3 extending Atari to 20M steps across three environments.
Each experiment produced per-environment results and cross-environment
aggregates using the rliable evaluation framework described in
Section~\ref{sec:eval-framework}.


\section{Synthesis of Results}

\subsection{Experiment 1: Classic-Control Benchmark}

On CartPole-v1, Acrobot-v1, and LunarLander-v3, PPO and A2C decisively
outperform the DQN family. PPO achieves a normalised IQM of 0.981 across
the three environments, with probabilities of improvement over every
other algorithm of at least 0.89. The DQN family (DQN and QR-DQN) fails
to learn at the allowed training budget (200K steps for CartPole and
Acrobot, 500K for LunarLander), scoring at or near random across all
seeds. RPPO collapses on Acrobot in a way that suggests its recurrent
policy is ill-suited to fully-observable MDPs. Classic-control is not a
DQN-friendly setting at these budgets.

\subsection{Experiment 2: Atari Benchmark}

Moving to Atari with CNN policies at 5M environment steps per seed
partially flips the per-environment ranking. On Pong, PPO continues to
dominate (mean return $\approx 20$ vs DQN $\approx -8$). On Breakout,
DQN slightly outperforms PPO, though both remain well below 11\% of the
theoretical maximum. The cross-environment aggregate still favours PPO
($\text{IQM}_{\text{PPO}} = 0.529$ vs $\text{IQM}_{\text{DQN}} = 0.130$,
$P(\text{PPO} > \text{DQN}) = 0.648$), but the margin is much smaller
than in Experiment~1, and the $M=2$ aggregate is visibly dominated by
Pong's dynamic range.

\subsection{Experiment 3: Extended Atari Benchmark}

Extending the budget to 20M steps and adding Seaquest as a third Atari
environment produces the most interesting result of the thesis: the
cross-environment aggregate \emph{flips}. The IQM now favours DQN
($0.443$ vs $0.277$), driven by DQN's $3.5\times$ Breakout advantage at
the larger budget. The pairwise probability of improvement, however,
still favours PPO ($P(\text{DQN} > \text{PPO}) = 0.393$). Per-environment,
DQN wins Breakout decisively (229.07 vs 65.50), PPO wins Pong (19.78 vs
11.20) and Seaquest (1173.92 vs 714.00) moderately.

The practical takeaway is that the IQM/POI disagreement is itself a
finding: with $M=3$ environments and 10 seeds, the choice of
cross-environment aggregation method can flip the answer to ``which
algorithm is better''. Experiment~3 also established that the training
pipeline is bitwise deterministic on the target RTX~4090 hardware --- a
practical result for anyone attempting to reproduce this work.


\section{Answer to the Research Question}


Neither family dominates. The correct comparison is per-environment,
and even within Atari the ranking depends on the environment, the
training budget, and the aggregation method:

\begin{itemize}
  \item On dense-reward, low-dimensional, fully-observable classic
  control, PPO and A2C are clearly stronger than DQN and QR-DQN at the
  budgets tested.
  \item On Atari Pong, PPO is consistently stronger at both 5M and 20M
  step budgets.
  \item On Atari Breakout, DQN edges PPO at 5M and wins decisively at
  20M.
  \item On Atari Seaquest at 20M, PPO wins moderately.
  \item Cross-environment aggregates are fragile at the $M$ values used
  here, and the choice of IQM vs pairwise POI can flip the aggregate
  answer. Per-environment reporting is the primary evidence.
\end{itemize}

The honest summary is that PPO is the stronger default --- it performs
well across all settings tested and dominates completely on classic
control --- but DQN closes the gap at longer Atari budgets and
outperforms PPO outright on Breakout. Both algorithms exist for good
reason, and the choice should be made based on the specific environment
and budget, not convention.


{\color{red} \section{Threats to Validity}}


\begin{itemize}
  \item \textbf{Default hyperparameters throughout.} No per-algorithm
  or per-environment tuning was performed. This keeps the comparison
  fair in the sense that no algorithm was given an advantage, but it
  means the numbers reflect ``out-of-the-box SB3'' rather than ``best
  achievable''. The gap between the two may be large, especially for
  DQN on Atari.
  \item \textbf{Small number of environments ($M$).} Classic control
  uses $M=3$, Atari uses $M=2$ (Experiment~2) then $M=3$
  (Experiment~3). These are well below the $M=50+$ used in standard RL
  benchmark papers. The aggregate metrics are correspondingly fragile:
  adding or removing one environment can change the cross-env aggregate
  substantially.
  \item \textbf{Fixed training budgets.} 200K--500K steps on classic
  control, 5M on Atari-Exp-2, 20M on Atari-Exp-3. These are smaller
  than the original DQN and PPO paper budgets (50M frames) and almost
  certainly disadvantage DQN more than PPO, since DQN's replay buffer
  requires more experience to warm up.
  \item \textbf{Experiment~3 buffer size oversight.} The intended DQN
  replay buffer size was 400K, but the experiments ran with 100K
  (Stable-Baselines3 default). The results are internally consistent
  and the determinism finding confirms exact reproducibility, but the
  Breakout lead attributed to DQN may look different with the intended
  buffer.
  \item \textbf{Timing data incomplete for Experiment~3 PPO.}
  Evaluation scores are complete for all seeds; only per-seed
  wall-clock timing has gaps for the PPO Pong runs. The timing section
  in Chapter~\ref{chap:exp3} uses a placeholder pending completion of
  the retraining.
\end{itemize}

